\section{结论}
本文围绕有限采样条件下的竞争源辨识，形成以双侧镜像嗅觉为主机制、判别式双机测量及持续反证为辅助机制的源概率定位研究稿。方法将母文中的方向判别算子转为联合源条件证据，以动态边界方向记忆约束测量提议，并在完整合法源集上进行成本分配与新后缀审计，最终输出二维源位置概率图。三个模块的创新归属、实际计算、信息权限及与普通联合观测和信息增益方法的区别均已给出。

已有原生归档对照表明，当前源对覆盖目标在仅有一个源对时提前饱和，动态镜像通道未被选入；在两个已标定源、4个独立开发世界中，普通单点和完整双点方法均4/4严格判对，M1为2/4，尚无独立定位收益。该结果将后续问题具体化为训练判别目标、有限编码和留出泛化之间的匹配，而不支持把所有失败统一归因于三维传播缺失、双机几何或镜像机制无效。

条件性理论说明了确定性编码的信息上限、成本分配的等价形式及正确观测族下的新后缀风险控制，并不构成真实羽流中的性能保证。下一阶段需先验证M1在同资源强基线之上的判别或样本效率收益，再开展双机主动几何、错误假设恢复和湖岸迁移。本文保留现有失败结果，将完整方法设计与实际验证状态分开，为后续独立实验提供明确可检验的命题。

\paragraph{研究稿及辅助写作说明}
本稿为2026年10月23日的V8研究讨论版本，未作为投稿定稿。使用ChatGPT（本会话模型标识GPT-6 Astra Pro）辅助原稿重构、文献与代码核查、公式整理、示意图排版及LaTeX编译；具体过程和核验边界列于配套审查文件。实验数据与成绩均来自已存在的冻结记录，未生成或改动原始观测。核心科研判断、方法与推导仍需研究者独立复核并对最终稿负责。《自动化学报》当前须知不仅要求披露辅助使用，还规定核心内容不应由AI工具生成\ccite{AASGuide}；采用其排版结构不表示本研究稿已满足其投稿条件，最终期刊和合规修改另行确认。
